\section{Introduction}

% Context
The integration of renewable energy sources into the power grid has driven the development of Virtual Power Plants (VPPs) to aggregate and manage distributed energy resources. Accurate forecasting of VPP dynamics---specifically electrical load, photovoltaic (PV) generation, and wind power generation---is critical for grid stability, economic dispatch, and operational risk management \cite{Ma2025,Cao2023,Qiu2024,Moreno2020,Wu2024}. Unlike single-variable forecasting, cooperative VPP management requires simultaneous, coherent predictions across these interrelated domains, where the physical constraints of each energy modality must be rigorously respected. For instance, PV generation is intrinsically bound by solar irradiance and cannot produce power during nighttime, while wind generation depends on a non-linear turbine power curve driven by localized wind speed. Reliable dispatch decisions hinge not only on minimizing raw statistical forecasting errors but also on ensuring that predictive models fundamentally comply with these immutable physical laws.

% Problem Statement
Recent advancements in deep learning have introduced powerful time-series forecasting architectures, yet their application to VPP scheduling presents fundamental challenges \cite{Vle2025,Arabzadeh2025}. While deep neural networks excel at minimizing statistical error metrics, they operate as black-box predictors that remain oblivious to the underlying physical principles governing power systems. In the present study's benchmark evaluations, the strongest Transformer baseline achieves a boundary violation rate (BVR) of 19.53\%, while recurrent architectures---despite their sequential inductive bias---still exhibit violation rates of 11.25\%. Consequently, pure data-driven frameworks frequently generate dimensionally accurate but physically impossible predictions, such as non-zero solar power during the night or negative wind generation, severely undermining operational safety and operator trust. Furthermore, the opaque nature of these architectures limits interpretability, making it exceedingly difficult to causally attribute anomalous predictions to specific meteorological inputs. Bridging the gap between the high predictive capacity of deep learning and the rigorous physical compliance required by coupled power systems remains a critical challenge. While heuristic post-processing---such as hard zero-clipping at physical boundaries---can retroactively enforce elementary static compliance (i.e., reducing BVR to zero), it strictly functions as a superficial post-hoc patch. Heuristic clipping inherently disrupts end-to-end continuous optimization by introducing non-differentiable dead zones, severing the natural continuous manifold of the underlying predictions. Such deterministic truncation fails to address the root cause of the model's structural ignorance, leaving the internal latent representations physically unbounded. This necessitates a framework that natively integrates differentiable physical priors upstream of the forecasting output, anchoring the neural network to intrinsic structural compliance.

% Contributions
This paper proposes PhysFormer, a physics-guided causal Transformer designed to occupy a definitively superior position on the accuracy-interpretability-physical compliance Pareto frontier for VPP forecasting. The proposed architecture fundamentally breaks the dichotomy between predictive acuity and structural rigor: it concurrently satisfies criteria across three critical dimensions previously unachieved in tandem. In the first dimension of accuracy, the model achieves a mean squared error (MSE) of 0.0128, highly competitive with the strongest black-box baseline. In the second dimension of interpretability, the framework achieves an orthogonal causal alignment that yields a deterministic gate-irradiance Pearson correlation coefficient of 0.8306, which remains robust ($r > 0.81$) even under severe 50\% exogenous sensor noise. In the third dimension of physical compliance, PhysFormer mathematically guarantees absolute structural safety, eliminating static boundary violations (BVR=0.00\%) through intrinsic manifold bounds. While pure data-driven models may achieve slightly lower temporal fluctuation rates through unconstrained over-fitting, PhysFormer treats temporal consistency as an engineering trade-off, prioritizing rigorous adherence to physical boundaries (BVR=0.00\%) over cosmetic smoothness.

The primary contributions of this paper are threefold. First, a Bounded Physical Act-Residual (BPAR) architecture is proposed. Fusing parallel Transformer and ExplicitPhysicalMapping streams via Softplus lower-bound clamping and physics activity masks, this mechanism intrinsically limits the neural residual search space. It mathematically guarantees that forecasting output never breaches normalized zero bounds, eradicating the boundary violation rate to an absolute 0.00\% without resorting to heuristic post-processing. Second, an orthogonal Physics-Guided Causal Coupling (PGCC) mechanism is introduced. By explicitly regularizing the cross-attention irradiance threshold learning, PGCC successfully imprints a highly correlated ($r=0.8306$) physical causality onto the temporal gating variable. Ablation studies reveal that incorporating this rigid physical footprint incurs negligible statistical accuracy loss ($\Delta\text{MSE}=0.0\%$), successfully decoupling causal attribution transparency from error minimization capacity while conferring extreme resilience to out-of-distribution (OOD) sensor noise. Third, the Ramp Violation Metric (RVM) is redefined within a \textit{safety-first} framework. We demonstrate that while unconstrained models like iTransformer exhibit exceptional smoothness, this is often a byproduct of ignoring physical limits (e.g., crossing zero-power thresholds during night). In contrast, PhysFormer's minor structural corrections at boundary limits (RVR-S=0.198\%) represent correct physical decisions to halt generation at the boundary, ensuring that absolute safety is prioritized over blind continuity. Specifically, PhysFormer guarantees BVR=0.00\% while bounding structural boundary corrections to physically negligible magnitudes.

% Paper Organization
The remainder of this paper is organized as follows. Section~\ref{sec:related} reviews related work in VPP forecasting, physics-informed neural networks, interpretability in time series, and physical metric design. Section~\ref{sec:method} details the proposed PhysFormer architecture, including the dual-stream design, physics activity gating, and physics-guided causal coupling mechanisms. Section~\ref{sec:training} describes the physics-constrained curriculum training strategy and loss function design. Section~\ref{sec:experiments} presents comparative experimental results, ablation studies, and physical interpretability analyses. Finally, Section~\ref{sec:conclusion} concludes the paper and discusses potential directions for future research.
